\documentclass[11pt]{article}
\usepackage[margin=1in]{geometry}
\usepackage{amsmath,amssymb,amsthm,mathtools}
\usepackage[T1]{fontenc}
\usepackage[utf8]{inputenc}
\usepackage{microtype}
\newtheorem{theorem}{Theorem}
\newtheorem{lemma}{Lemma}
\newtheorem{remark}{Remark}
\newcommand{\R}{\mathbb R}
\newcommand{\C}{\mathbb C}
\newcommand{\boxt}{\mathbin{\boxtimes}}

\title{A log-uniform counterexample to free multiplicative log-unimodality}
\author{ProofNet24 closure note for Hasebe--Ueda--Wang Problem 4.3}
\date{16 September 2026}

\begin{document}
\maketitle

\begin{abstract}
Hasebe, Ueda and Wang asked whether the free positive multiplicative convolution
of two multiplicatively symmetric log-unimodal probability measures must again
be log-unimodal.  We give an explicit negative answer.  Let
\[
 d\mu(x)=\frac{1}{12x}{\bf 1}_{[e^{-6},e^6]}(x)\,dx.
\]
Then $\log_*\mu$ is uniform on $[-6,6]$, but if $\rho=\mu\boxt\mu$ and $g$
is the density of $\log_*\rho$, then $g''(0)>0$.  Thus the symmetry point is a
strict local minimum and $\rho$ is not log-unimodal.  The proof is an explicit
$\eta$/$\Sigma$-transform calculation with a certified rational curvature bound.
\end{abstract}

\section{Source question and the example}

For a probability measure $\lambda$ on $\R_+=(0,\infty)$, Hasebe--Ueda--Wang
call $\lambda$ log-unimodal when $\log_*\lambda$ is unimodal on $\R$.
Multiplicative symmetry means invariance under $x\mapsto x^{-1}$.
Their Problem 4.3 asks whether
\[
 \mu,\nu \text{ multiplicatively symmetric and log-unimodal}
 \quad\Longrightarrow\quad
 \mu\boxt\nu \text{ log-unimodal}.
\]
This is different from their theorem for convolution with the free positive
multiplicative Brownian law $\sigma_t$.

Set $L=6$, $a=e^{-L}$, and
\begin{equation}\label{eq:mu}
 d\mu_L(x)=\frac{1}{2Lx}{\bf 1}_{[e^{-L},e^L]}(x)\,dx.
\end{equation}
Under $x=e^u$, the measure in \eqref{eq:mu} becomes $du/(2L)$ on $[-L,L]$.
Hence $\mu_L$ is multiplicatively symmetric and log-unimodal.  We prove that
$\rho_L:=\mu_L\boxt\mu_L$ is not log-unimodal.

\section{Exact inverse transform}

For a probability measure $\lambda$ on $\R_+$ write
\[
 \psi_\lambda(z)=\int_{\R_+}\frac{zx}{1-zx}\,d\lambda(x),
 \qquad
 \eta_\lambda(z)=\frac{\psi_\lambda(z)}{1+\psi_\lambda(z)}.
\]
For \eqref{eq:mu}, direct integration gives
\begin{equation}\label{eq:psi}
 \psi_L(z)=\frac{1}{2L}\log\!\left(\frac{1-az}{1-a^{-1}z}\right).
\end{equation}
If $w=\eta_L(z)$, then $\psi_L(z)=w/(1-w)$.  Put
\[
 E(w)=\exp\!\left(\frac{2Lw}{1-w}\right).
\]
Exponentiating \eqref{eq:psi} and solving for $z$ yields the exact inverse
\begin{equation}\label{eq:q}
 q_L(w):=\eta_L^{-1}(w)
   =\frac{1-E(w)}{a-a^{-1}E(w)}.
\end{equation}
Since $\Sigma_\lambda(w)=\eta_\lambda^{-1}(w)/w$ and $\Sigma$-transforms
multiply under free positive multiplicative convolution,
\begin{equation}\label{eq:Q}
 Q_L(w):=\eta_{\rho_L}^{-1}(w)=\frac{q_L(w)^2}{w}.
\end{equation}

The inverse branch in \eqref{eq:q} is the physical one.  Indeed, on the
negative real axis
\[
 \psi_\lambda'(z)=\int_{\R_+}\frac{x}{(1-zx)^2}\,d\lambda(x)>0,
 \qquad z<0,
\]
so $\eta_\lambda$ is strictly increasing there and its inverse is unique.
At $w=-1$ one has $E(-1)=a$, hence $q_L(-1)=-1$ and $Q_L(-1)=-1$.
This anchors the physical inverse branch at $w=-1$.

\section{A noncritical unit-circle continuation}

Write
\[
 w=e^{i\theta}=e^{2it},\qquad 0<t<\frac\pi2.
\]
Then
\[
 \frac{w}{1-w}=-\frac12+\frac{i}{2}\cot t,
 \qquad
 E(w)=a e^{i\phi},
 \qquad
 \phi=L\cot t.
\]
Let
\[
 B_a(v)=\frac{v-a}{1-av}.
\]
Formula \eqref{eq:q} becomes
\[
 q_L(w)=-\frac{1}{B_a(e^{i\phi})}.
\]
Thus $|q_L(w)|=|Q_L(w)|=1$.  Choose the continuous argument
$B_a(e^{i\phi})=e^{i\beta_a(\phi)}$ for $0\leq\phi\leq\pi$.
The derivative is the Poisson kernel
\[
 \beta_a'(\phi)=P_a(\phi)
 =\frac{1-a^2}{1-2a\cos\phi+a^2}.
\]
Consequently, along the unit-circle arc,
\begin{equation}\label{eq:phase}
 Q_L(e^{i\theta})=e^{i\varphi_L(\theta)},
 \qquad
 \varphi_L(\theta)=2\pi-\theta-2\beta_a\!\left(L\cot\frac\theta2\right),
\end{equation}
and
\begin{equation}\label{eq:phaseprime}
 \varphi_L'(\theta)
 =L\csc^2 t\,P_a(\phi)-1.
\end{equation}

We need only elementary explicit bounds.  Since
\[
 e>\sum_{k=0}^5\frac1{k!}=\frac{163}{60}
 \quad\text{and}\quad
 \left(\frac{163}{60}\right)^6>400,
\]
we have $a=e^{-6}<1/400$.  Therefore
\[
 P_a(\phi)\geq\frac{1-a}{1+a}>\frac{399}{401},
\]
and \eqref{eq:phaseprime} gives the uniform lower bound
\begin{equation}\label{eq:Apositive}
 \varphi_L'(\theta)>
 6\frac{399}{401}-1=\frac{1993}{401}>0.
\end{equation}

Let $t_0=\arctan(L/\pi)$, so that $L\cot t_0=\pi$.  At
$\theta_0=2t_0$ we have $B_a(-1)=-1$ and the continuous phase in
\eqref{eq:phase} equals $-2t_0<0$.  At $\theta=\pi$ it equals $\pi$.
By \eqref{eq:Apositive}, there is a unique
\[
 \theta_*=2t_*\in(2t_0,\pi)
\]
for which $\varphi_L(\theta_*)=0$, hence
\begin{equation}\label{eq:centerroot}
 Q_L(w_*)=1,\qquad w_*=e^{2it_*}.
\end{equation}
Moreover $Q_L'$ never vanishes along the arc from $-1$ to $w_*$.  Analytic
continuation from the negative-axis inverse therefore identifies this entire
arc, in particular $w_*$, with the physical inverse branch.  No algebraic
spurious branch enters the calculation below.

\section{The center curvature}

Let $g_L$ denote the density of the logarithmic push-forward
$\log_*\rho_L$.  Near $y=0$, write the physical boundary inverse as $w=w(y)$,
so that
\begin{equation}\label{eq:boundaryeq}
 Q_L(w(y))=e^{-y},\qquad w(0)=w_*.
\end{equation}
The elementary relation between the Cauchy transform and $\eta$ is
\[
 G_\rho(\zeta)=\frac{\zeta^{-1}}{1-\eta_\rho(\zeta^{-1})}.
\]
Taking the Stieltjes boundary value and using the upper-half-plane conjugate
branch gives
\begin{equation}\label{eq:density}
 g_L(y)=\frac1\pi\Im\frac{1}{1-w(y)}.
\end{equation}

Set
\[
 A=\varphi_L'(2t_*)>0,
 \qquad
 B=\varphi_L''(2t_*).
\]
Differentiating \eqref{eq:boundaryeq}, equivalently differentiating the phase
relation first and then inverting it, gives
\[
 w'(0)=-\frac{w_*}{A},
 \qquad
 w''(0)=\frac{w_*(A+iB)}{A^3}.
\]
Using
\[
 \frac{w_*}{(1-w_*)^2}=-\frac14\csc^2t_*,
\]
\[
 \frac{w_*^2}{(1-w_*)^3}
 =\frac18\csc^2t_*
  -\frac{i}{8}\csc^2t_*\cot t_*,
\]
in the second derivative of \eqref{eq:density} yields
\begin{equation}\label{eq:gsecond}
 g_L''(0)
 =-\frac{\csc^2t_*}{4\pi A^3}
   \bigl[B+A\cot t_*\bigr].
\end{equation}

It remains to determine the sign of the bracket.  Write
\[
 D=1-2a\cos\phi+a^2,
 \qquad \phi=L\cot t_*,
\]
where here and below $\phi$ is evaluated at $t_*$.  Differentiating
\eqref{eq:phaseprime} and simplifying gives the exact cancellation
\begin{equation}\label{eq:cancellation}
 B+A\cot t_*
 =-\cot t_*
 +\frac{L^2a(1-a^2)\sin\phi}{D^2\sin^4t_*}.
\end{equation}
Because $t_*>t_0$, we have $0<\phi<\pi$.  Hence
$\sin\phi<\phi=L\cot t_*$.  Also
\[
 D\geq(1-a)^2,
 \qquad
 \sin t_*>\frac{L}{\sqrt{L^2+\pi^2}}.
\]
Substitution into \eqref{eq:cancellation} gives
\begin{equation}\label{eq:CL}
 B+A\cot t_*<\cot t_*\,[-1+C_L],
 \qquad
 C_L=\frac{a(1-a^2)(L^2+\pi^2)^2}{L(1-a)^4}.
\end{equation}

Now use $\pi^2<(22/7)^2<10$ and define
\[
 F(L)=\frac{e^{-L}(L^2+10)^2}{L(1-e^{-L})^4}.
\]
Then $C_L<F(L)$ and, for $L\geq6$,
\[
 \frac{d}{dL}\log F(L)
 =-1+\frac{4L}{L^2+10}-\frac1L
   -\frac{4e^{-L}}{1-e^{-L}}
 <-1+\frac3L\leq-\frac12.
\]
Thus $F$ is decreasing on $[6,\infty)$.  At $L=6$, using $e^{-6}<1/400$,
\begin{equation}\label{eq:rational}
 F(6)
 <\frac1{400}\frac{46^2}{6}\left(\frac{400}{399}\right)^4
 =\frac{67712000000}{76034875203}<1.
\end{equation}
Therefore $C_6<1$, and \eqref{eq:CL} shows
\[
 B+A\cot t_*<0.
\]
By \eqref{eq:gsecond},
\begin{equation}\label{eq:strictmin}
 \boxed{g_6''(0)>0.}
\end{equation}

\section{Conclusion}

The input $\mu_6$ is invariant under $x\mapsto x^{-1}$.  Inversion commutes
with free positive multiplicative convolution at the level of distributions,
so $\rho_6=\mu_6\boxt\mu_6$ is also multiplicatively symmetric and $g_6$ is
even.  Equation \eqref{eq:strictmin} says that its symmetry point is a strict
local minimum.  An even unimodal density cannot have a strict local minimum at
its center.  We have proved the following.

\begin{theorem}
Let
\[
 d\mu(x)=\frac{1}{12x}{\bf 1}_{[e^{-6},e^6]}(x)\,dx.
\]
Then $\mu$ is multiplicatively symmetric and log-unimodal, but
$\mu\boxt\mu$ is not log-unimodal.  Hence Hasebe--Ueda--Wang Problem 4.3 has a
negative answer.
\end{theorem}

\begin{remark}
The same estimates prove the stronger family statement for every $L\geq6$.
The fixed value $L=6$ is presented here because one counterexample completely
answers the source question and is the common scope of the two independent
end-to-end audits in the accompanying closure package.
\end{remark}

\begin{remark}[Finite certificate and novelty scope]
The exact rational inequalities used above were executed in the canonical
ProofNet24 artifact `log\_uniform\_L6\_family\_exact\_check\_v2.py`; the archived
output is `PASS`.  Two independent OTHER-track audits rederived the source-level
argument.  A bounded literature search through 16 September 2026 did not locate
a published resolution matching this counterexample, but this note makes no
publication-grade novelty claim on the strength of a negative search.
\end{remark}

\begin{thebibliography}{9}
\bibitem{HUW}
T. Hasebe, Y. Ueda and J.-C. Wang,
\emph{Log-unimodality for free positive multiplicative Brownian motion},
arXiv:2009.13848 (2020); Colloq. Math. 169 (2022), 209--226.
\end{thebibliography}

\end{document}
